% Paper 5, §4 — Contamination: the boundary between "well-known" and "leaked"
% Anchors: kb/notes/evaluation/eval-methodology.md §contamination
%          kb/notes/evaluation/knowledge-benchmarks.md §contamination
%          kb/excerpts/{contamination-survey-2025,mmlu-redux-2024}.md

\section{Contamination: the boundary between well-known and leaked}
\label{sec:contamination}

The \glsterm{mmlu}{MMLU} lineage of \cref{sec:knowledge-benchmarks} closes on the
question this section opens. A frontier model that scores within the
verifier-error band of a saturated knowledge benchmark either
\emph{knows} the material in the deployment-relevant sense or has
\emph{seen} the material during pre-training in a way that lets it
reproduce the answer \glsterm{key}{key}. The boundary between those two states is
fuzzy by construction --- the same exam-bank items that make \glsterm{mmlu}{MMLU} a
broad-coverage benchmark are exactly the items that pre-training
scrapes from public web text --- and the field has not converged on a
protocol that draws the boundary cleanly. The thesis of this section
is that contamination has three structurally distinct failure modes
(direct, indirect, methodological), four detection methods with
characteristic false-positive / false-negative profiles, one
worked-case-study where the boundary is genuinely contested
(\textsc{aime~2024}), and zero field-standard contamination protocols
shipped alongside benchmark releases. The next layer of the layered
defense (\cref{sec:reasoning-benchmarks}) inherits all of this, because
reasoning benchmarks are knowledge benchmarks with longer derivation
chains and narrower answer distributions.

\subsection{Three failure modes}
\label{sec:contamination-modes}

\citet{contamination-survey-2025} structures the field's response into
\emph{static} and \emph{dynamic} evaluation strategies%
\footnote{KB excerpt: \texttt{kb/excerpts/contamination-survey-2025.md\#abstract}.}
\citep[\S{}abstract]{contamination-survey-2025}:

\begin{quote}\itshape
To mitigate the risk of potential \glsterm{data-contamination}{data contamination}, \glsterm{llm}{LLM} benchmarking
has undergone a transformation from static to dynamic benchmarking.
\ldots\ We first examine methods that enhance static benchmarks and
identify their inherent limitations. We then highlight a critical
gap--the lack of standardized criteria for evaluating dynamic
benchmarks.
\end{quote}

The taxonomy below is the synthesis the contamination notes carry
forward%
\footnote{KB note: \texttt{kb/notes/evaluation/eval-methodology.md\#sec-2-2}; the canonical category list in the survey itself is \texttt{[PDF-VERIFY]}-flagged. \deepencite{contamination-survey-2025 §sec-taxonomy}}
into three failure modes that differ in where the leak enters the
pipeline.

\paragraph{Direct contamination.} A test-set item appears verbatim or
near-verbatim in the model's pre-training corpus. The benchmark
question $\rho$ and (often) the gold answer $a^\star$ are tokens in the
training sequence; the model's score on that item is partially a
recall of the training distribution rather than a measurement of the
construct the benchmark advertises. This is the failure mode that
$n$-gram-overlap \glsterm{decontamination}{decontamination} filters target. \citet{hendrycks2021-mmlu}'s
\glsterm{mmlu}{MMLU} was constructed from publicly-available exam banks (AP exams,
GRE, MCAT, professional licensing exams), which is exactly the
distribution that web crawls preserve at high redundancy.

\paragraph{Indirect contamination.} The verbatim item is absent, but a
paraphrase, translation, partial overlap, distractor-set match, or
worked solution is present. \citet{contamination-survey-2025}'s
abstract names static-benchmark enhancement as one of the survey's
two structural concerns, with the implicit recognition that the
enhancements (paraphrase, distractor swap, language transfer) only
shift contamination from \emph{input} to \emph{label} or
\emph{distribution}: a model trained on AIME 2023 problems has seen
the AIME style, the AIME difficulty calibration, and the AIME answer
distribution, and that exposure transfers to AIME 2024 even when no
2024 item is in the training set%
\footnote{KB note: \texttt{kb/notes/evaluation/eval-methodology.md\#sec-2-2} enumerates input/label/distribution/indirect/generative variants.}.
The StackExchange-and-Reddit explanations that surround any widely-used
benchmark are themselves an indirect contamination channel that no
$n$-gram filter on the benchmark text can catch.

\paragraph{Methodological contamination.} The leak happens at
\emph{evaluation time}, not training time. The model is given
tool-use, retrieval, or fine-tune-on-eval access during scoring, so
the gold answer (or material that determines it) re-enters the inference
context. This is structurally different from the first two: no
amount of training-corpus auditing closes it, because the channel is
the eval harness itself. \textsc{rag}-augmented model cards that
report \textsc{mmlu} scores under tool-use, agentic-benchmark
scaffolds that browse the live web, and \textsc{\glsterm{situational-awareness}{llm}}-as-judge
pipelines whose judge has been fine-tuned on the same task are all
instances. The contamination-survey abstract does not name this mode
explicitly, but the dynamic-benchmark gap it calls out --- ``the lack
of standardized criteria for evaluating dynamic benchmarks''
\citep[\S{}abstract]{contamination-survey-2025} --- is partly the
methodological-contamination problem dressed in a positive frame.

\begin{intuition}
The three modes correspond to three distinct repair points in the
layered-defense framing of \cref{sec:introduction}. \emph{Direct}
contamination is repaired upstream of training, by \glsterm{decontamination}{decontamination}
filters on the corpus. \emph{Indirect} contamination is repaired at
benchmark construction, by contamination-resistant items (Google-proof
\cref{sec:knowledge-gpqa}; unpublished \cref{sec:knowledge-frontiermath}).
\emph{Methodological} contamination is repaired downstream, in the
eval harness's isolation guarantees. A benchmark that defends against
only one of the three is not contaminated-clean; it is
contaminated-clean-along-one-axis. (Returning to canonical form: the
score random variable $\hat{S}(\theta, \mathcal{B}, \pi, \tau)$ from
\cref{sec:eval-methodology} carries an additive bias term whose
support is the union of the three contamination channels, and zeroing
out one channel does not zero the bias.)
\end{intuition}

\subsection{Four detection methods, four error profiles}
\label{sec:contamination-detection}

No detection method is uniformly best; each has a characteristic
false-positive / false-negative profile that determines what the
method can rule in and what it can rule out. The four methods below
are the field's working set as of 2026%
\footnote{KB note: \texttt{kb/notes/evaluation/eval-methodology.md\#sec-5-1} catalogs these methods and their adoption status.}.

\paragraph{$n$-gram overlap.} Flag any training document that contains
a benchmark item's $n$-gram of length $\geq k$ (Brown et al.\ used
$k = 13$ tokens for GPT-3 \deepencite{brown2020-gpt3 §C-decontamination}).
\textbf{Cheap, standard, the only method most pre-training pipelines
actually run.} False-positive profile: low; a 13-gram match is rarely
coincidental. False-negative profile: severe under any paraphrase,
translation, or partial-template rewrite. The method catches direct
contamination and misses essentially all indirect contamination.

\paragraph{Embedding-space neighbors.} Embed both training documents
and benchmark items in a shared dense space and flag training docs
whose nearest-neighbor distance to a benchmark item is below threshold.
\textbf{Catches paraphrases the $n$-gram filter misses.} False-positive
profile: high and threshold-dependent --- legitimate exam-bank-style
items not in the benchmark cluster near benchmark items by genre, not
by content. False-negative profile: depends on the embedding model's
own training distribution; a benchmark item paraphrased into a domain
the embedding model under-represents will be missed.

\paragraph{Prompted recall (Min-K\% Prob, perplexity gap).} For a
candidate item, compute a model-internal memorization signal:
$\sum_{i \in K\%\text{-lowest}} \log P(x_i)$ \deepencite{shi2024-min-k};
or the ratio of the model's \glsterm{perplexity}{perplexity} on the held-out test split
to its \glsterm{perplexity}{perplexity} on training-distribution-matched text. Memorized
items have characteristic flat-low-\glsterm{perplexity}{perplexity} signatures.
\textbf{Model-internal --- does not need access to the training corpus.}
False-positive profile: confounded by item difficulty and stylistic
familiarity; an easy item the model would always assign high
probability to looks contaminated even when it is not. False-negative
profile: degraded by post-training (\glsterm{rlhf}{RLHF} / instruction-tuning) which
overwrites base-model probability surfaces; the signal is strongest
on base models and weakens substantially on chat-tuned variants.

\paragraph{Behavioral probes.} Construct items the contaminated model
should fail on if it is reasoning rather than memorizing: shuffle the
multiple-choice option order, swap the gold-letter position,
substitute synonymous distractors, perturb numerical inputs.
A model whose accuracy collapses on the perturbed variant is leaning
on memorized surface form. \textbf{Falsifiable in either direction.}
False-positive profile: a model that legitimately handles the
canonical surface but is brittle to perturbation looks contaminated
when it is merely brittle. False-negative profile: a model that
memorized the underlying \emph{logic} (not the surface form) of a
contaminated item passes the perturbation test even though the
training distribution leaked the answer.

The methods are non-redundant: each catches a different leak channel,
and the correlation between their flags across a benchmark is far
from one. The field's open-source \glsterm{decontamination}{decontamination} pipelines run a
subset (typically just $n$-gram overlap); vendor pipelines reportedly
run more (LLaMA 3 and Phi-4 model cards document multi-method
filters), but the audits are vendor-internal and not externally
reproducible.

\subsection{The \textsc{aime~2024} case study}
\label{sec:contamination-aime}

\textsc{aime~2024} is the worked example. The American Invitational
Mathematics Examination's 2024 problems were administered in early
2024, then publicly released alongside official answer keys, and have
since served as the canonical \emph{reasoning} benchmark for late-2024
and 2025-frontier models. \citet{deepseek-r1}'s headline result is
the load-bearing illustration%
\footnote{KB excerpt: \texttt{kb/excerpts/deepseek-r1.md\#sec-2-3-aime}.}
\citep[\S{}2.3]{deepseek-r1}:

\begin{quote}\itshape
Figure 1(a) depicts the performance trajectory of DeepSeek-\glsterm{r1-zero}{R1-Zero} on
the AIME 2024 benchmark throughout the RL training process, where the
average pass@1 score on AIME 2024 shows a significant increase,
jumping from an initial 15.6\% to 77.9\%.
\end{quote}

The headline jump from $15.6\%$ to $77.9\%$ at $\mathrm{pass}@1$
(further to $86.7\%$ under \glsterm{self-consistency-2}{self-consistency}) is the load-bearing
RL-on-verifiable-rewards demonstration of Paper~3 (§reasoning-RL).
What the contamination question asks of that result is whether the RL
optimization is teaching new reasoning computation or learning to
invoke pretrained computation that already saw enough \textsc{aime}-style
material to do this --- the open question that the
reasoning-training KB note flags as the load-bearing 2026 dispute%
\footnote{KB note: \texttt{kb/notes/reasoning/reasoning-training.md\#sec-3-2}.}.

The timeline that complicates the question:

\begin{enumerate}
\item \textbf{February 2024:} \textsc{aime~2024} administered;
problems and official answer keys publicly released within weeks.
\item \textbf{Q2--Q3 2024:} \textsc{aime~2024} problems and worked
solutions appear on Art-of-Problem-Solving forums, math-coaching
blogs, and StackExchange threads --- the indirect-contamination
channel of \cref{sec:contamination-modes}.
\item \textbf{Q4 2024:} pre-training data cutoffs for late-2024
frontier models include the public-web representation of these
problems and (for problems with worked solutions) their
derivation chains. Reasoning-RL traces are subsequently generated
on these same problems, sometimes from a base model that has
already seen them.
\item \textbf{2025:} \textsc{aime~2024} is the canonical
reasoning-RL headline benchmark, used to compare model classes
that have differing levels of training-time exposure to its
items.
\end{enumerate}

The line between ``publicly canonical'' (the problems are by
construction released for human study) and ``contaminated'' (the
items are in the model's training distribution in a way the
benchmark's measurement assumes they are not) is, on
\textsc{aime~2024}, structurally fuzzy. The same release that makes
the benchmark legible to human contestants makes it absorbable by
crawl-based pre-training. The dispute is not over a methodological
mistake; it is over which benchmark category \textsc{aime~2024} is
in. Paper~3's discussion of the same headline picks up this
contradiction explicitly (\hyperref[P3-sec:contradictions]{Paper~3 §14 (AIME-contamination)}).

\begin{contradiction}
\textbf{What counts as contamination?} The vendor-side position is
that \emph{decontaminated}-by-the-released-procedure equals clean:
LLaMA 3, Phi-4, and other frontier model cards document
$n$-gram-overlap-based filters and report decontaminated training
sets, and on that basis report \textsc{aime}-2024-style scores as
measurements of reasoning capability. The third-party position is
that the survey of leak channels in \cref{sec:contamination-modes}
makes this position incomplete: $n$-gram filters miss
indirect/generative contamination, contamination-survey
\citep[\S{}abstract]{contamination-survey-2025} explicitly calls out
``the lack of standardized criteria'' as a gap, and the
multi-method audits that would close the loop are not run
field-wide. As of 2026-Q2, no contamination protocol ships alongside
benchmark releases the way \texttt{requirements.txt} ships alongside
Python code; each lab certifies its own pipeline against its own bar.
The two positions are not formally inconsistent --- \glsterm{decontamination}{decontamination}
\emph{by some procedure} is a real artifact --- but they make
incompatible claims about what reported benchmark scores mean,
because they implicitly fix different procedures as the procedure.
\end{contradiction}

\subsection{Why the field has not standardized}
\label{sec:contamination-no-standard}

The mechanical reason is that the four detection methods of
\cref{sec:contamination-detection} have non-overlapping cost-quality
profiles, and no single subset is dominated by another. $n$-gram
overlap is cheap and reproducible but indirectly-blind;
embedding-neighbor methods catch indirect leakage but require
expensive training-corpus access and tunable thresholds;
prompted-recall methods are model-internal and corpus-free but
weakened by post-training; behavioral probes are cheap and
falsifiable but confounded by brittleness. A field-standard protocol
would have to specify which subset to run, at what threshold, on
which corpus slice --- and the right answer plausibly differs across
benchmarks. \textsc{mmlu-redux} \citep{mmlu-redux-2024} is a partial
analogue at the verifier-error layer (a published per-question
correction set that subsequent runs can use), and its partial
adoption --- the Open \glsterm{llm}{LLM} Leaderboard moved its primary headline to
\textsc{mmlu-pro} while many model cards still report raw \textsc{mmlu}
\citep[\S{}implications]{mmlu-redux-2024} --- is the diagnostic for
how field-wide adoption of a heavier contamination protocol would
likely play out.

The structural reason is that the incentives diverge. A vendor
reporting on its own model is rewarded for a clean score and for a
documented \glsterm{decontamination}{decontamination} story; a third party auditing the same
score has no comparable reward. \citet{contamination-survey-2025}'s
EMNLP-2025 framing is the closest the field has come to a
field-standard \emph{discussion}, but the survey is a description of
what exists, not a protocol that ships. The methodological
dynamic-benchmark proposals (\glsterm{matharena}{MathArena}, LiveCodeBench, GAIA's gated
leaderboard) work around contamination rather than measure it; that
is the right move where it is feasible, but it abandons the question
of what to do with the legacy benchmarks the field continues to
report on.

\subsection{Forward-link}
\label{sec:contamination-forward}

The next layer of the layered defense, in
\cref{sec:reasoning-benchmarks}, asks whether benchmarks designed for
multi-step reasoning escape the contamination problem this section
named. The short answer is that they inherit it: \textsc{aime} and
\textsc{math-500} are exactly the case study above; \textsc{gpqa
diamond} \citep{rein2023-gpqa} mitigates direct contamination by
\emph{construction} but cannot rule out indirect signal from worked
PhD-level expositions on the open web; \textsc{frontiermath}
\citep{glazer2024-frontiermath} is the most aggressive defense in the
lineage with its unpublished-item construction, at the cost of being
single-source and not externally replicable. The contradiction this
section opened --- ``what counts as contamination?'' --- is
collected with its sibling contradictions in
\cref{sec:contradictions}, where the closing-condition for each is
named: a field-standard contamination-detection protocol shipped
alongside benchmark releases, the way reproducibility tooling now
ships alongside published scientific code.
